\chapter*{Περίληψη}

Η Ανάκτηση-Επαυξημένη Παραγωγή (Retrieval-Augmented Generation, RAG) έχει αναδειχθεί
ως ένα θεμελιώδες παράδειγμα για την αγκύρωση των απαντήσεων των Μεγάλων Γλωσσικών
Μοντέλων (ΜΓΜ) σε εξωτερικές, επαληθεύσιμες πηγές γνώσης. Παρά την αξιόλογη πρόοδο
που έχει επιτευχθεί σε μονόστροφα σενάρια RAG, το πολύστροφο συνομιλιακό πλαίσιο
— όπου το σύστημα πρέπει να παράγει ακριβείς και αξιόπιστες απαντήσεις στο πλαίσιο
εκτεταμένου διαλόγου — παραμένει σημαντικά δυσκολότερο και σχετικά αδύναμα μελετημένο.
Στις πολύστροφες αλληλεπιδράσεις, η ποιότητα ανάκτησης υποβαθμίζεται διαδοχικά λόγω
μη αυτοτελών ερωτημάτων, αναφορών αντωνυμιών και σιωπηρής μεταφοράς θέματος, ενώ
η ποιότητα παραγωγής επιδεινώνεται περαιτέρω από σωρευτικά σφάλματα, αναπάντητες
ερωτήσεις και την αυξανόμενη πολυπλοκότητα του συνομιλιακού πλαισίου.
 
Η παρούσα διπλωματική εργασία παρουσιάζει μια ενιαία, σταθεροποιημένη αρχιτεκτονική
για πολύστροφο RAG, η οποία αναπτύχθηκε και αξιολογήθηκε στο σύνολο δεδομένων
αναφοράς MTRAG στο πλαίσιο της κοινής εργασίας SemEval-2026 Task 8. Το σύστημα
αντιμετωπίζει και τις τρεις υποεργασίες: ανάκτηση αποσπασμάτων (Task A),
παραγωγή απαντήσεων με δοθέντα αποσπάσματα αναφοράς (Task B) και ολοκληρωμένο
RAG από άκρο σε άκρο (Task C). Η κεντρική αρχή σχεδιασμού είναι η σταθερότητα
ανάκτησης μέσω ποικιλότητας ερωτημάτων: αντί να βασιστούμε σε ετερογενή σύνολα
ανακτητών, εκδίδουμε πέντε συμπληρωματικές αναδιατυπώσεις ερωτημάτων βασισμένες
σε ΜΓΜ — Ελάχιστη, Corpus-Ειδική, Υποθετική Ενσωμάτωση Εγγράφου (HyDE), Αλυσίδα
Σκέψης (CoT) και Λέξεις-Κλειδιά Άγκυρας — σε έναν μοναδικό αραιό ανακτητή
ευθυγραμμισμένο με το σώμα κειμένων (ELSER v1), και συγκεντρώνουμε τις κατατεταγμένες
λίστες μέσω ενός ιεραρχικού σχήματος Εμφωλευμένης Αντίστροφης Σύντηξης Κατατάξεων
(Nested Reciprocal Rank Fusion, RRF).
 
Για την παραγωγή απαντήσεων, υιοθετούμε έναν πρακτορο-κεντρικό αγωγό πολλαπλών
σταδίων που αποσυνθέτει την αγκυρωμένη παραγωγή σε πέντε διαδοχικά στάδια:
προκαταρκτικό έλεγχο απαντησιμότητας, εξαγωγή αποδεικτικών εκτάσεων, δημιουργία δύο
υποψήφιων απαντήσεων, επιλογή μέσω κριτών και τελική μικροδιόρθωση. Μια συνάρτηση
διαμόρφωσης εκχύλισης 4-γραμμάτων επιβάλλει συμβιβασμό μεταξύ πιστότητας και
φυσικότητας, τιμωρώντας τόσο απαντήσεις επιρρεπείς σε ψευδαισθήσεις (hallucination)
όσο και υπερβολικά κυριολεκτικές αντιγραφές. Ο αγωγός ολοκληρωμένου RAG ενσωματώνει
επιπλέον μια πύλη απαντησιμότητας πολλαπλών κριτών με ψηφοφορία σταθμισμένη ως
προς την εμπιστοσύνη.
 
Μελέτες αφαίρεσης στοιχείων στο development set δείχνουν ότι: (i) η ποικιλότητα ερωτημάτων
πάνω σε έναν μοναδικό ανακτητή υπερτερεί της σύνθεσης ετερογενών ανακτητών, με βάση
τις επίσημες κρίσεις συνάφειας του benchmark· (ii) η εξαγωγή αποδεικτικών εκτάσεων πριν
την παραγωγή βελτιώνει την πιστότητα· (iii) η βαθμονόμηση της απαντησιμότητας γίνεται ο καθοριστικός παράγοντας της ολοκληρωμένης
απόδοσης όταν αυξάνεται το ποσοστό των μη απαντήσιμων ερωτημάτων, ενώ τα σφάλματα ανάκτησης
στις κορυφαίες θέσεις παραμένουν σημαντική δεύτερη πηγή αποτυχιών. Το σύστημά μας
κατατάχθηκε 1ο στο Task A (nDCG@5: 0.5776, +20.5\% έναντι της ισχυρότερης γραμμής
βάσης), 2ο στο Task B (Αρμονικός Μέσος: 0.7698) και 11ο στο Task C (Αρμονικός Μέσος:
0.5409) στον επίσημο πίνακα κατάταξης SemEval-2026, μεταξύ 38, 26 και 29 συμμετοχών
αντίστοιχα.
 
\paragraph*{Λέξεις-κλειδιά ---}
Ανάκτηση-Επαυξημένη Παραγωγή, Πολύστροφη Συνομιλία, Αναδιατύπωση Ερωτημάτων,
Αντίστροφη Σύντηξη Κατατάξεων, Πρακτορο-Κεντρική Παραγωγή, Ανίχνευση
Απαντησιμότητας, Μεγάλα Γλωσσικά Μοντέλα, Ανάκτηση Πληροφοριών, SemEval-2026.
